\documentclass[a4paper,12pt,twocolumn,landscape]{article}

\usepackage{FabZ}
\usepackage{vecteurs}
\usepackage{repere}

\usepackage{geometry}
\geometry{hmargin=0.5cm,vmargin=1.5cm}

\newcommand{\milieu}[6]{$\pointcoord{A_{\theenumi}}{#1}{#2}$ et $\pointcoord{B_{\theenumi}}{#3}{#4}$ \hfill \reponseEX{$\pointcoord{M_{\theenumi}}{#5}{#6}$}}

\newcommand{\extremite}[6]{$\pointcoord{A_{\theenumi}}{#1}{#2}$ et $\pointcoord{M_{\theenumi}}{#3}{#4}$ \hfill \reponseEX{$\pointcoord{B_{\theenumi}}{#5}{#6}$}}

\newcommand{\longueur}[5]{$\pointcoord{A_{\theenumi}}{#1}{#2}$ et $\pointcoord{B_{\theenumi}}{#3}{#4}$ \hfill \reponseEX{Longueur $A_{\theenumi}B_{\theenumi}$ = $#5$}}

\newcommand{\quadrilatere}[9]{$\pointcoord{A_{\theenumi}}{#1}{#2}$, $\pointcoord{B_\theenumi}{#3}{#4}$ $\pointcoord{C_{\theenumi}}{#5}{#6}$, $\pointcoord{D_\theenumi}{#7}{#8}$ \\ \textcolor{red}{$A_{\theenumi}B_{\theenumi}C_{\theenumi}D_{\theenumi}$ est un #9.}}

\newcommand{\parallelogrammeManqueA}[8]{$\pointcoord{B_{\theenumi}}{#1}{#2}$, $\pointcoord{C_{\theenumi}}{#3}{#4}$ $\pointcoord{D_{\theenumi}}{#5}{#6}$ \hfill \textcolor{red}{$\pointcoord{A_{\theenumi}}{#7}{#8}$}}

\newcommand{\parallelogrammeManqueB}[8]{$\pointcoord{A_{\theenumi}}{#1}{#2}$, $\pointcoord{C_{\theenumi}}{#3}{#4}$ $\pointcoord{D_{\theenumi}}{#5}{#6}$ \hfill \textcolor{red}{$\pointcoord{B_{\theenumi}}{#7}{#8}$}}

\newcommand{\parallelogrammeManqueC}[8]{$\pointcoord{A_{\theenumi}}{#1}{#2}$, $\pointcoord{B_{\theenumi}}{#3}{#4}$ $\pointcoord{D_{\theenumi}}{#5}{#6}$ \hfill \textcolor{red}{$\pointcoord{C_{\theenumi}}{#7}{#8}$}}

\newcommand{\parallelogrammeManqueD}[8]{$\pointcoord{A_{\theenumi}}{#1}{#2}$, $\pointcoord{B_{\theenumi}}{#3}{#4}$ $\pointcoord{C_{\theenumi}}{#5}{#6}$ \hfill \textcolor{red}{$\pointcoord{D_{\theenumi}}{#7}{#8}$}}

%
%\setlength{\headheight}{0pt}
%\pagestyle{fancyplain}
%\fancyhf{}
%\lhead[]{\textbf{TD Vecteurs}}
%\chead[]{}
%\rhead[]{}
%
%\lfoot[]{}
%\cfoot[]{}
%\rfoot[]{}
%%\rfoot[]{Page \thepage~ sur \pageref{LastPage}}

\fancypagestyle{firststyle}
{
	\setlength{\headheight}{0em}
	\fancyhf{}
	\lhead[]{\textbf{Mathématiques\\DS2}}
	\chead[]{Seconde 12\\ lundi 12 octobre 2015}
	\rhead[]{Séquence 1\\ Repérage dans le plan}

	\lfoot[]{}
	\cfoot[]{}
	\rfoot[]{}
%	\rfoot[]{Page \thepage~ sur \pageref{LastPage}}
}

\fancypagestyle{vide}
{
	\setlength{\headheight}{0em}
	\pagestyle{fancyplain}
	\def\headrulewidth{0em}
	\fancyhf{}
	\lhead[]{}
	\chead[]{}
	\rhead[]{}
	
	\lfoot[]{}
	\cfoot[]{}
	\rfoot[]{}
%	\rfoot[]{Page \thepage~ sur \pageref{LastPage}}
}

\usetikzlibrary{fadings}
\newcommand{\Fin}{node[xshift=-1.5ex,rotate=10]{F}
node[rotate=170]{i}
node[xshift=1.5ex,rotate=45]{n}}

\newcommand{\bonnesvacances}{node[xshift=0ex,rotate=10]{B}
node[xshift=1*1.5ex,rotate=170]{o}
node[xshift=2*1.5ex,rotate=45]{n}
node[xshift=3*1.5ex,rotate=128]{n}
node[xshift=4*1.5ex,rotate=75]{e}
node[xshift=5*1.5ex,rotate=130]{s}
node[xshift=6*1.5ex,rotate=85]{~}
node[xshift=7*1.5ex,rotate=128]{v}
node[xshift=8*1.5ex,rotate=43]{a}
node[xshift=9*1.5ex,rotate=4]{c}
node[xshift=10*1.5ex,rotate=145]{a}
node[xshift=11*1.5ex,rotate=5]{n}
node[xshift=12*1.5ex,rotate=25]{c}
node[xshift=13*1.5ex,rotate=105]{e}
node[xshift=14*1.5ex,rotate=45]{s}
}


\begin{document}

\vspace*{1em}

\begin{minipage}{0.45\textwidth}
\thispagestyle{firststyle}

\vspace*{1em}

\exercice\hfill \emph{4 points}

Un élève a voulu tracer un parallélogramme $BDCA$ mais il ne sait pas où placer le sommet $A$.\\ Comment pouvez-vous l'aider~?\\ Vous détaillerez votre démarche sur votre copie avant de réaliser des calculs.

\begin{center}

\begin{tikzpicture}[scale=1,every node/.style={scale=1}]

%\coordinate (X) at (-0.5,-0.5);
%\coordinate (Y) at (13,13);
%\clip (X) rectangle (Y);
%\placerpoint{A}{-4}{1}{below left};
\placerpoint{B}{5}{3}{above right};
\placerpoint{C}{0}{4}{below right};
\placerpoint{D}{9}{6}{above left};
%\placerpoint{M}{6}{4}{above left};
%\draw[very thick,red] (A) -- (B) -- (C) -- (D) -- cycle;
\repereOIJ{-1}{10}{-1}{7};
\foreach \r in {5}
    \draw[thick, below right] (\r,0) node{\r};
\foreach \r in {5}
    \draw[thick, above left] (0,\r) node{\r};

\end{tikzpicture}

\end{center}


\vspace*{2em}

\exercice\hfill \emph{6 points}

On se place dans le repère~\emph{orthonormé}\rep{O}{I}{J}. Dans ce repère, on considère les points $\pointcoord{R}{-2}{0}$, $\pointcoord{S}{0}{2}$ et $\pointcoord{T}{7}{-7}$.\\Vos réponses seront argumentées.

	\begin{enumerate}
		\item Le point $R$ appartient-il au cercle $\mathscr{C}$ de centre $T$ passant par~$S$~?
		\item Le point $T$ appartient-il à la médiatrice de $\left[RS\right]$~?
%		\item Déterminer les coordonnées de $Z$ pour que $AZBC$ soit un parallélogramme.
	\end{enumerate}



\vspace{-2em}


\end{minipage}
\newpage

\vspace*{1em}

\begin{minipage}{0.45\textwidth}
\thispagestyle{firststyle}

\vspace*{1em}

%\paragraph{Exercice~4} 
%
%On se place dans le repère \emph{orthonormé}\rep{O}{I}{J}. $ABCD$ est un rectangle de centre $O$. $A$ est tel que son abscisse est négative et son ordonnée positive. Le côté $AB$ est parallèle à l'axe des ordonnées. Pour chaque question, faire une figure en prenant $AB = 4 cm$ et $AD = 6 cm$. Donner les coordonnées des points $A$, $B$, $C$, $D$, $O$, $I$ et $J$~:
%	\begin{enumerate}
%		\item Dans le repère\rep{O}{I}{J}
%		\item Dans le repère\rep{B}{C}{A}
%		\item Dans le repère\rep{A}{B}{D}
%	\end{enumerate}

%\vspace{-2em}

\exercice\hfill \emph{10 points}

Le plan $\left(\mathscr{P}\right)$ est muni d'un repère~\emph{orthonormé}\rep{O}{I}{J}.\\ Dans ce repère, on considère les points suivants $\pointcoord{A}{-10}{-11}$ ; $\pointcoord{B}{2}{-2}$, $\pointcoord{C}{11}{10}$ et $\pointcoord{D}{-1}{1}$.
On désigne par $E$, $F$ et $G$ les milieux respectifs des segments $\left[AB\right]$,  $\left[BC\right]$ et $\left[AC\right]$.

\begin{enumerate}
	\item Démontrez que le triangle $BCD$ est isocèle en $D$.
	\item Démontrez que le quadrilatère $ABCD$ est un losange.
	\item Calculez les coordonnées des points $E$, $F$ et $G$ dans le repère\rep{O}{I}{J}.
	\item Quelle est la nature du triangle $EFG$~? Justifiez.
	\item Déterminez les coordonnées des points $A$, $B$, $C$, $D$, $E$, $F$ et $G$ dans le~repère\rep{A}{B}{C}.
	\item Pourquoi ne peut-on pas calculer des distances dans le~repère\rep{A}{B}{C}~?
\end{enumerate}

\vspace*{1em}

\vspace*{\stretch{1}}
\begin{center}
\begin{tikzpicture}[scale=2,transform shape]
\draw (0,0) \Fin;
\draw (-1em,-1ex) -- (1em,-1ex);
\path[scope fading=south] (-1em,-0.25em) rectangle (1em,-3.75ex);
\draw[yscale=-1] (0,2ex) \Fin;
\end{tikzpicture}
\end{center}
\vspace*{\stretch{1}}

\vspace*{2em}

\hrulefill
\exercice\hfill \emph{Bonus}

Les diagonales d’un quadrilatère $ABCD$ se coupent en $E$.

$I$, $J$, $K$, $L$ sont les milieux respectifs de $[AB]$, $[BC]$, $[CD]$, $[DA]$.

\begin{enumerate}
	\item Faire une figure en y reportant toutes les informations de l’énoncé.
	\item Démontrer que $IJKL$ est un parallélogramme.
\end{enumerate}


\end{minipage}
\end{document}
